% !TeX spellcheck = en_US
\documentclass[french]{yLectureNote}

\title{Outils Mathématiques 2}
\subtitle{Physique}
\author{Paulhenry Saux}
\date{\today}
\yLanguage{Français}

\professor{F.Pettinari}
\usepackage{graphicx}%----pour mettre des images
\usepackage[utf8]{inputenc}%---encodage
\usepackage{geometry}%---pour modifier les tailles et mettre a4paper
%\usepackage{awesomebox}%---pour les boites d'exercices, de pbq et de croquis ---d\'esactiv\'e pour les TP de PC
\usepackage{tikz}%---pour deiffner + d\'ependance de chemfig
% \usepackage{tabularx}%---pour dimensionner automatiquement les tableaux avec variable X
\usepackage{awesomebox}%---Pour les boites info, danger et autres
\usepackage{menukeys}%---Pour deiffner les touches de Calculatrice
\usepackage{fancyhdr}%---pour les en-t\^ete personnalis\'ees
\usepackage{blindtext}%---pour les liens
\usepackage{hyperref}%---pour les liens (\`a mettre en dernier)
\usepackage{caption}%---pour la francisation de la l\'egende table vers Tableau
\usepackage{pifont}
\usepackage{array}%---pour les tableaux
\usepackage{lipsum}
\usepackage{yFlatTable}
\usepackage{multicol}
\newcommand{\Lim}[1]{\lim\limits_{\substack{#1}}\:}
\renewcommand{\vec}{\overrightarrow}
\newcommand{\N}[0]{\mathbb{N}}
\newcommand{\dd}{\mathrm{d}}
\newcommand{\norm}[1]{||\vec{#1}||}
\newcommand{\fo}{\psi(\vec{r},t)}
\newcommand{\foe}{\psi(\vec{r},t)\*}
\newcommand{\HH}{\hat{H}}
\newcommand{\hb}{\hbar}
\newcommand{\lap}{\nabla^2}
\newcommand{\lapcc}{\frac{\partial^2 }{\partial x^2}+\frac{\partial^2 }{\partial y^2}+\frac{\partial^2 }{\partial z^2}}
\newcommand{\mpsi}{\(\psi\)}
\begin{document}
\setcounter{chapter}{3}
\chapter{Séries numériques }
\section{Suites et séries numériques}
\subsection{Séries numérique}
\begin{definition}[Séries numériques]
\[S_n = \sum_{n=n_0}^N u_n\]
\end{definition}
\begin{definition}[Série géométrique]
Soit \(S_N = \sum_{n=0}^N q^n\). Si \(q\neq 1\), on a \(S_N=\frac{1-q^{N+1}}{1-q}\). On en déduit que la série converge si \(|q|<1\)
\end{definition}
\begin{definition}[Série de Riemann]
Une série de la forme \(\frac{1}{n^{\alpha}}\) est une série de Riemann. Elle converge si \(\alpha>1\)
\end{definition}
\begin{proposition}[divergence grossière]
Si la série de terme général \(U_n\) converge, alors \(\lim u_n =0\)
\end{proposition}
\begin{proposition}[Critère des séries alternées]
Pour une série alternée de terme \(u_n\), si la suite \(|u_n|\) est décroissante et tend vers 0, alors la série converge.
\end{proposition}
\begin{proposition}[Convergence absolue]
Si une série convege absolument (la suite de terme \(|u_n|\) converge), elle converge.
\end{proposition}
\begin{proposition}[Critère de d'Alembert]
Si \(\frac{|u_{n+1}|}{|u_n|}\to l\). Alors si \(l<1\) la série converge, si \(l>1\), elle diverge.
\end{proposition}
\subsection{Séries de fonctions}
\begin{theorem}
 Si il y a convergence uniforme, alors on peut inverser les limites : \[\Lim{x\to x_0}\Lim{n\to\infty}f_n(x) = \Lim{n\to\infty}\Lim{x\to x_0}f_n(x)\]
\end{theorem}
\section{Séries entières}
\begin{definition}[Série entière]
C'est la suite de polynomes \(a_nx^n\)
\end{definition}
\begin{definition}[Rayon de convergence]
Noté R c'est une valeur telle que \(\forall x_0, |x_0|<R\) la série entière converge aboslument.
\end{definition}
\begin{theorem}[Règle de D'Alembert]
 On suppose \(a_n\neq 0\) pour n suffisament grand. Si \(\Lim{n\to\infty}\big|\frac{a_{n+1}}{a_n}\big|\to l\), alors \(R=\frac{1}{l}\)
\end{theorem}
\subsection{Séries à connaitre}
\begin{theorem}[DL de l'exponentielle]
 \[e^x = \sum_{n=0}^{+\infty} \frac{x^n}{n!}, \quad R = +\infty\]
\end{theorem}
\begin{theorem}[DL de cosinus]
 \[\cos(x) = \sum_{n=0}^{+\infty} (-1)^n \frac{x^{2n}}{(2n)!}, \quad R = +\infty\]
\end{theorem}
\begin{theorem}[DL de sinus]
 \[\sin(x) = \sum_{n=0}^{+\infty} (-1)^n \frac{x^{2n+1}}{(2n+1)!}, \quad R = +\infty\]
\end{theorem}
\begin{theorem}[DL de cosinus hyperbolique]
 \[\cosh(x) = \sum_{n=0}^{+\infty} \frac{x^{2n}}{(2n)!}, \quad R = +\infty\]
\end{theorem}
\begin{theorem}[DL de sinus hyperbolique]
 \[\sinh(x) = \sum_{n=0}^{+\infty} \frac{x^{2n+1}}{(2n+1)!}, \quad R = +\infty\]
\end{theorem}
\begin{theorem}[Série géométrique]
 \[\frac{1}{1+x} = \sum_{n=0}^{+\infty} (-1)^n x^n, \quad R = 1\]
\end{theorem}
\begin{theorem}[DL de ln]
 \[\ln(1+x) = \sum_{n=1}^{+\infty} (-1)^{n-1} \frac{x^n}{n}, \quad R = 1\]
\end{theorem}
\begin{theorem}[DL de la série géométrique au carré]
 \[\frac{1}{(1-x)^2} = \sum_{n=0}^{+\infty} (n+1)x^n, \quad R = 1\]
\end{theorem}
\begin{theorem}[DL de la série géométrique à la puissance quelconque]
 \[\frac{1}{(1-x)^p} = \sum_{n=0}^{+\infty} \binom{n+p-1}{p-1} x^n, \quad R = 1\]
\end{theorem}
\end{document}
